\documentclass[a4paper,11pt]{article}
\usepackage[francais]{babel}
%% Prévu pour compiler avec lualatex
 \usepackage[utf8]{inputenc}
%\usepackage{fontspec}
%\usepackage{libertine}
\usepackage[T1]{fontenc}
\usepackage{graphicx}
\usepackage{fancyhdr}
\usepackage[top=2.5cm, bottom=2.5cm, left=2cm, right=2cm]{geometry}
\usepackage{listings}
%\usepackage[utf8]{luainputenc}
\usepackage[hidelinks]{hyperref}
\usepackage{caption}

\lfoot{\bsc{Enseirb-Matmeca}}
\rfoot{Informatique --- 3\ieme{} année}

\pagestyle{fancy}
\begin{document}

\begin{titlepage}
  \begin{center}

    \begin{center}
      \includegraphics[width=4cm]{EM.jpg}
    \end{center}

    \vspace*{1cm}
        
    \rule{0.75\linewidth}{0.7mm}\\[0.4cm]
    {\Huge Rapport TP6 --- Lancer de rayons\\[0.4cm]}
    \rule{0.75\linewidth}{0.7mm} \\[1.5cm]

    {\Large Bazire \bsc{Houssin}\\Sylvain \bsc{Vaglica}\\Stéphane \bsc{Castelli}\\[2cm]}
    {\Large Dimanche 2 Février 2014}
  \end{center}
\end{titlepage}

\tableofcontents
\clearpage
\section{Introduction}

Afin de représenter une scène, potentiellement dynamique, sur un ordinateur, une technique courante est le lancer de rayons (\emph{ray tracing}). Celle-ci se base sur le principe de retour inverse de la lumière, en effet au lieu de calculer les trajectoires des rayons de lumière pour chacune des sources, ce qui représente un gaspillage de ressources car la plupart de ceux-ci n'atteignent pas l'observateur, on calcule la trajectoire de rayons de lumière qui partiraient de l'observateur pour aller vers les surfaces environnantes. Ainsi, on ne calcule que ce qui est observable. Cette technique a ceci de particulier qu'elle est fortement parallélisable, car chacun des rayons est indépendant des autres. En partant de la version séquentielle de l'algorithme, on implémentera une version parallèle. Ensuite, on mettera en pratique la technique du \emph{vol de travail}.

\section{Version parallèle}



\section{Conclusion}

Réaliser une factorisation LU en utilisant MPI nécessite en premier lieu de bien comprendre comment la réaliser. En effet, comprendre ses mécanismes s'est effectué essentiellement par la lecture de code en Fortran, il a donc fallu assimiler la méthode --- ce qui a demandé un certain temps --- en conséquence de quoi la parallélisation, pourtant cœur du projet, n'a pas pu être aussi bien avancée qu'il l'aurait été souhaité. Toutefois le principe a pu être compris, et des progrès considérables auraient pu être accomplis en un laps de temps minime. En outre, cela nous a permis  d'approfondir notre capacité à mettre en pratique l'algèbre numérique couplée à la programmation informatique.

\end{document}
